\chapter{Conclusion} \label{Chap6}\label{chap:conclusion}

\section{Research outcome}

This dissertation asked to what extent overlapping satellite-conditioned
predictions can be combined into a coloured regional mesh and integrated with
geospatial context in ChordAtlas.  The study establishes both parts of this
workflow.  A rectangular large-image method fuses density predictions at
sub-voxel positions, replaces abrupt overlap changes with smooth weights, and
recovers colour on the fixed global surface.  The resulting mesh is then
placed in a shared spatial frame, associated with OSM footprints and DSM
context, separated into modular assets, and published to ChordAtlas.

The matched London experiment showed substantially lower window-boundary
discontinuity without increasing the measured GPU-memory peak, and the colour
pass covered the complete surface without changing its geometry.  The current
on-demand integration retained that colour through coordinate conversion, DSM
processing, footprint extraction, Java parsing, and building publication.
Manifests, stable selection identities, cache checks, and explicit stage states
make the process inspectable across the Python--Java boundary.

\section{Principal contributions}

The work makes three main contributions:

\begin{enumerate}
  \item It extends large-image Sat3DGen inference to rectangular mosaics through
  sub-voxel density placement and smooth overlap fusion, and recovers global
  vertex colour in a memory-bounded second pass without altering the extracted
  geometry.
  \item It defines a coordinate-explicit path from fused or independently
  assembled geometry to footprint-associated ChordAtlas assets, preserving
  selection identity and vertex appearance while adding OSM and DSM context.
  \item It introduces and evaluates manifest, cache, validation, and
  publication contracts across Python and Java, using data audits, matched
  fusion measurements, coordinate checks, topology inspection, and controlled
  regression cases.
\end{enumerate}

These contributions provide a practical connection between learned urban
reconstruction and procedural authoring.  Sat3DGen supplies site-conditioned
density and appearance, while OSM, DSM, and ChordAtlas add spatial identity,
height context, modular asset boundaries, and an environment for further
editing.  CityEngine clarifies the complementary value of rule-driven
semantics and controllability, while the retained ChordAtlas exports illustrate
the representation difference between per-vertex colour and structured
procedural material maps.

\section{Implications and next steps}

The evaluation identifies the next research priorities.  Multi-area tests with
independent geometry and appearance references are needed to determine whether
smoother window transitions also improve reconstruction accuracy.  The
independent-mesh path requires corrected vertical stitching and more selective
DSM processing, while the published building assets require topology repair.
For deployment, density and colour should be bound to one immutable model
revision, the large-image scripts should be packaged with the bridge, and
readiness should include Java conversion and visible loading acknowledgement.

The current CityEngine export and retained ChordAtlas representation examples
can be extended into a controlled authoring study with shared identity, inputs,
and edit tasks, while
Ubiq can test interactive and collaborative use of the resulting assets.

\section{Closing statement}

The dissertation connects bounded satellite predictions to a continuous
coloured surface and an auditable urban-asset workflow.  Its central value lies
in combining learned
density and appearance with explicit coordinates, geographic identity, and
publication state, so that generated geometry can participate in a larger
editable system.  This provides a concrete foundation for evaluating regional
reconstruction, procedural authoring, and collaborative XR with the same
spatially identified assets.
